% Exemplo: matemática com tkz-euclide (triângulo com ângulos e circunferência inscrita), pgfplots
% (área sob a curva), cancel e mathtools (simplificação passo a passo), tabela-verdade e
% bussproofs (dedução natural).
% Compilar: latexmk -pdf matematica.tex           Testado com: tkz-euclide 5, pgfplots 1.18, TeX Live 2024 ou mais recente
%%BEGIN:matematicapreambulo
\documentclass[12pt, a4paper]{article}
\usepackage[T1]{fontenc}
\usepackage[brazilian]{babel}
\usepackage{mathtools, amssymb}  % mathtools carrega o amsmath
\usepackage{cancel}              % \cancel{...}: termo riscado
\usepackage{tkz-euclide}         % geometria euclidiana (carrega o TikZ)
\usepackage{pgfplots}            % gráficos de funções
\pgfplotsset{compat=1.18}
\usepgfplotslibrary{fillbetween} % pinta a região entre duas curvas
\usepackage{bussproofs}          % árvores de prova (dedução natural)
%%END:matematicapreambulo

\begin{document}
\pagestyle{empty}

%%BEGIN:matematicageometria
Triângulo retângulo em $C$, com a circunferência inscrita de centro $I$:
\begin{center}
\begin{tikzpicture}
  \tkzDefPoints{0/0/C, 4/0/A, 0/3/B}
  \tkzDefTriangleCenter[in](A,B,C) \tkzGetPoint{I}
  \tkzDefPointBy[projection=onto A--C](I) \tkzGetPoint{T}
  \tkzDrawPolygon(A,B,C)
  \tkzDrawCircle[gray](I,T)
  \tkzMarkRightAngle(A,C,B)
  \tkzMarkAngle[size=0.8](B,A,C) \tkzLabelAngle[pos=1.15](B,A,C){$\alpha$}
  \tkzDrawPoints(A,B,C,I)
  \tkzLabelPoints[below](A,C) \tkzLabelPoints[left](B)
  \tkzLabelPoints[above right](I)
  \tkzLabelSegment[below](C,A){$b = 4$}
  \tkzLabelSegment[left](C,B){$a = 3$}
  \tkzLabelSegment[above right](A,B){$c = 5$}
\end{tikzpicture}
\end{center}
%%END:matematicageometria

%%BEGIN:matematicaintegral
A área sob $f(x) = x^2/2 + 1$ entre $x = 1$ e $x = 3$ é
$\int_1^3 f(x)\,dx = \frac{19}{3}$:
\begin{center}
\begin{tikzpicture}
  \begin{axis}[axis lines=middle, xlabel={$x$}, ylabel={$y$},
      xmin=0, xmax=4, ymin=0, ymax=7, xtick={1,3}, ytick=\empty,
      width=7cm, height=5.5cm]
    \addplot[name path=f, thick, domain=0:3.6, samples=50]
      {x^2/2 + 1} node[left, pos=0.9] {$f$};
    \path[name path=eixo] (axis cs:0,0) -- (axis cs:4,0);
    \addplot[gray!30] fill between[of=f and eixo,
      soft clip={domain=1:3}];
  \end{axis}
\end{tikzpicture}
\end{center}
%%END:matematicaintegral

%%BEGIN:matematicasimplificacao
\begin{align*}
  \lim_{x \to 2} \frac{x^2 - 4}{x - 2}
    &= \lim_{x \to 2} \frac{(x + 2)\cancel{(x - 2)}}{\cancel{x - 2}}
    && \text{diferença de quadrados} \\
    &= \lim_{x \to 2} (x + 2) = 4
    && \text{substituindo } x = 2
\end{align*}
\begin{align*}
  (a + b)^2 - (a - b)^2
    &= \underbrace{a^2 + 2ab + b^2}_{(a + b)^2}
     - \underbrace{(a^2 - 2ab + b^2)}_{(a - b)^2} \\
    &= \cancel{a^2} + 2ab + \cancel{b^2} - \cancel{a^2} + 2ab
     - \cancel{b^2} = 4ab
\end{align*}
%%END:matematicasimplificacao

\clearpage
%%BEGIN:matematicalogica
\begin{center}
\begin{tabular}{cc|cccc}
  $p$ & $q$ & $p \land q$ & $p \lor q$ & $p \to q$
    & $p \leftrightarrow q$ \\ \hline
  V & V & V & V & V & V \\
  V & F & F & V & F & F \\
  F & V & F & V & V & F \\
  F & F & F & F & V & V \\
\end{tabular}
\end{center}
Conjuntos: $\mathbb{N} \subset \mathbb{Z} \subset \mathbb{Q}
\subset \mathbb{R}$, $A = \{x \in \mathbb{N} \mid x < 5\}$,
$A \cup B$, $A \cap B$, $A \setminus B$, $\overline{A}$
e $\emptyset$.
%%END:matematicalogica

%%BEGIN:matematicaprova
Dedução natural: \emph{modus ponens} e a comutatividade do $\land$.
\begin{center}
\AxiomC{$p \to q$}
\AxiomC{$p$}
\RightLabel{\small ${\to}E$}
\BinaryInfC{$q$}
\DisplayProof
\qquad
\AxiomC{$p \land q$}
\RightLabel{\small ${\land}E_2$}
\UnaryInfC{$q$}
\AxiomC{$p \land q$}
\RightLabel{\small ${\land}E_1$}
\UnaryInfC{$p$}
\RightLabel{\small ${\land}I$}
\BinaryInfC{$q \land p$}
\DisplayProof
\end{center}
%%END:matematicaprova

\end{document}
